\documentclass[10pt,a4paper]{amsart}
\usepackage{amsmath,amssymb,amsthm,graphicx}
\usepackage{verbatim}
\newenvironment{reference}{\comment}{\endcomment}
%\newenvironment{reference}{}{}
\newenvironment{slogan}{\comment}{\endcomment}
\newenvironment{history}{\comment}{\endcomment}
% For commutative diagrams we use Xy-pic
\usepackage[all]{xy}

% We use 2cell for 2-commutative diagrams.
\xyoption{2cell}
% Theorem environments.
%
\theoremstyle{plain}
\newtheorem{theorem}[subsection]{Theorem}
\newtheorem{proposition}[subsection]{Proposition}
\newtheorem{lemma}[subsection]{Lemma}

\theoremstyle{definition}
\newtheorem{definition}[subsection]{Definition}
\newtheorem{example}[subsection]{Example}
\newtheorem{exercise}[subsection]{Exercise}
\newtheorem{situation}[subsection]{Situation}

\theoremstyle{remark}
\newtheorem{remark}[subsection]{Remark}
\newtheorem{remarks}[subsection]{Remarks}

\numberwithin{equation}{subsection}

% Macros
%
\def\lim{\mathop{\mathrm{lim}}\nolimits}
\def\colim{\mathop{\mathrm{colim}}\nolimits}
\def\Spec{\mathop{\mathrm{Spec}}}
\def\Hom{\mathop{\mathrm{Hom}}\nolimits}
\def\Ext{\mathop{\mathrm{Ext}}\nolimits}
\def\SheafHom{\mathop{\mathcal{H}\!\mathit{om}}\nolimits}
\def\SheafExt{\mathop{\mathcal{E}\!\mathit{xt}}\nolimits}
\def\Sch{\mathit{Sch}}
\def\Mor{\mathop{\mathrm{Mor}}\nolimits}
\def\Ob{\mathop{\mathrm{Ob}}\nolimits}
\def\Sh{\mathop{\mathit{Sh}}\nolimits}
\def\NL{\mathop{N\!L}\nolimits}
\def\CH{\mathop{\mathrm{CH}}\nolimits}
\def\proetale{{pro\text{-}\acute{e}tale}}
\def\etale{{\acute{e}tale}}
\def\QCoh{\mathit{QCoh}}
\def\Ker{\mathop{\mathrm{Ker}}}
\def\Im{\mathop{\mathrm{Im}}}
\def\Coker{\mathop{\mathrm{Coker}}}
\def\Coim{\mathop{\mathrm{Coim}}}

% Boxtimes
%
\DeclareMathSymbol{\boxtimes}{\mathbin}{AMSa}{"02}

%
% Macros for moduli stacks/spaces
%
\def\QCohstack{\mathcal{QC}\!\mathit{oh}}
\def\Cohstack{\mathcal{C}\!\mathit{oh}}
\def\Spacesstack{\mathcal{S}\!\mathit{paces}}
\def\Quotfunctor{\mathrm{Quot}}
\def\Hilbfunctor{\mathrm{Hilb}}
\def\Curvesstack{\mathcal{C}\!\mathit{urves}}
\def\Polarizedstack{\mathcal{P}\!\mathit{olarized}}
\def\Complexesstack{\mathcal{C}\!\mathit{omplexes}}
% \Pic is the operator that assigns to X its picard group, usage \Pic(X)
% \Picardstack_{X/B} denotes the Picard stack of X over B
% \Picardfunctor_{X/B} denotes the Picard functor of X over B
\def\Pic{\mathop{\mathrm{Pic}}\nolimits}
\def\Picardstack{\mathcal{P}\!\mathit{ic}}
\def\Picardfunctor{\mathrm{Pic}}
\def\Deformationcategory{\mathcal{D}\!\mathit{ef}}
\usepackage{cleveref}
\crefname{originallemma}{original lemma}{original lemmas}
\Crefname{originallemma}{Original lemma}{Original lemmas}
\makeatletter
\newlabel{morphisms-lemma-differential-product}{{01V1}{source}{}{sourcepointer-1}{}}
\newlabel{morphisms-lemma-differential-product@cref}{{[originallemma][0][]01V1}{[1][0][]source}{}{}{}}
\makeatother
\title{The de Rham complex of a product is the total complex of the exterior box product of the factors}
\author{The Stacks Project Authors}
\date{stacks-project:derham.tex:local-snapshot-20261012; GNU FDL 1.2 or later}
\begin{document}
\maketitle
\noindent\textit{Editorial note.} All delivered mathematical text is the author's own, copied from the pinned Stacks Project chapter file derham.tex; no statement or proof step was rewritten, reordered or removed.\par\medskip
\noindent\textit{Editorial note.} Preamble provenance: the pinned chapter file opens with \textbackslash{}input\{preamble\} and carries no preamble text of its own, so the project's single Stacks-wide preamble staged at sources/stacks-derham/preamble.tex is excerpted verbatim as three spans, lines 8-12 (verbatim package and the project's own reference/slogan/history display environments), lines 14-18 (Xy-pic setup) and lines 154-217 (theorem declarations and mathematical macros). The wrapper class and the amsmath/amssymb/amsthm/graphicx base packages are supplied by the assembler. No macro was invented or redefined.\par\medskip
\noindent\textit{Editorial note.} Same-chapter pointers are located at their \textbackslash{}label lines inside the pinned chapter file; cross-chapter pointers are located in the project's own tag list staged in the snapshot (sources/stacks-chow/tags.txt), at the verbatim lines carrying their canonical tags.\par\medskip
\noindent\textit{Allowed standard prerequisites (editorial).}\begin{itemize}
\item The differentials of a product as a direct sum. Morphisms, label morphisms-lemma-differential-product (invoked at proof line 886). It identifies \$\textbackslash{}Omega\_\{X \textbackslash{}times\_S Y/S\}\$ with the direct sum \$p\textasciicircum{}*\textbackslash{}Omega\_\{X/S\} \textbackslash{}oplus q\textasciicircum{}*\textbackslash{}Omega\_\{Y/S\}\$, which is the input to the degreewise computation.
\item Elementary properties of exterior powers and the definition of the box product. Standard, used at proof lines 887-894. Taking exterior powers gives \$\textbackslash{}Omega\textasciicircum{}n\_\{X \textbackslash{}times\_S Y/S\} = \textbackslash{}bigoplus\_\{i+j=n\} \textbackslash{}Omega\textasciicircum{}i\_\{X/S\} \textbackslash{}boxtimes \textbackslash{}Omega\textasciicircum{}j\_\{Y/S\}\$, matching the terms of the total complex.
\end{itemize}\medskip
\ifdefined\hypertarget\hypertarget{sourcepointer-1}{}\fi
\noindent\textit{Source reference.} Original lemma 01V1: sources/stacks-chow/tags.txt, source lines 2430--2430. Navigation pointer for \textbackslash{}ref\{morphisms-lemma-differential-product\} used inside the counted span (line 886); the argument that consumes it is retained in the span.\par\medskip
\begin{lemma}
\label{lemma-de-rham-complex-product}
In the situation above there is a canonical isomorphism
$$
\text{Tot}(\Omega^\bullet_{X/S} \boxtimes \Omega^\bullet_{Y/S})
\longrightarrow
\Omega^\bullet_{X \times_S Y/S}
$$
of complexes of $f^{-1}\mathcal{O}_S$-modules.
\end{lemma}

\begin{proof}
We know that
$
\Omega_{X \times_S Y/S} = p^*\Omega_{X/S} \oplus q^*\Omega_{Y/S}
$
by Morphisms, Lemma \ref{morphisms-lemma-differential-product}.
Taking exterior powers we obtain
$$
\Omega^n_{X \times_S Y/S} =
\bigoplus\nolimits_{i + j = n}
p^*\Omega^i_{X/S} \otimes_{\mathcal{O}_{X \times_S Y}} q^*\Omega^j_{Y/S} =
\bigoplus\nolimits_{i + j = n}
\Omega^i_{X/S} \boxtimes \Omega^j_{Y/S}
$$
by elementary properties of exterior powers. These identifications determine
isomorphisms between the terms of the complexes on the left and the right
of the arrow in the lemma. We omit the verification that these maps
are compatible with differentials.
\end{proof}

\end{document}
